%% kapa: exam=1 round=33 year=2022 period=1 form=A subject=real_estate subject_ko=부동산학원론 seq=10 number=90
%% source: Q-Net 감정평가사 기출문제 (cst00309) artlSeq=5209997 file=raw/1008/5209997/2022년도 제33회 감정평가사 1차시험 1교시(민법,경제학,부동산학).pdf
%% conversion: kordoc 4.24.51 -> tools/gen1.py (auto); review=proofread-sonnet; flags=none
\begin{question}{90}
주거분리와 여과과정에 관한 설명으로 옳은 것은?
\begin{choices}
  \item 여과과정이 원활하게 작동하면 신규주택에 대한 정부지원으로 모든 소득계층이 이득을 볼 수 있다.
  \item 하향여과는 고소득층 주거지역에서 주택의 개량을 통한 가치상승분이 주택개량비용보다 큰 경우에 발생한다.
  \item 다른 조건이 동일할 경우 고가주택에 가까이 위치한 저가주택에는 부(-)의 외부효과가 발생한다.
  \item 민간주택시장에서 불량주택이 발생하는 것은 시장실패를 의미한다.
  \item 주거분리현상은 도시지역에서만 발생하고, 도시와 지리적으로 인접한 근린지역에서는 발생하지 않는다.
\end{choices}
\end{question}
\answer{1}
